% Options for packages loaded elsewhere
\PassOptionsToPackage{unicode}{hyperref}
\PassOptionsToPackage{hyphens}{url}
\documentclass[
  ukrainian,
]{article}
\usepackage{xcolor}
\usepackage[margin=2.5cm]{geometry}
\usepackage{amsmath,amssymb}
\setcounter{secnumdepth}{-\maxdimen} % remove section numbering
\usepackage{iftex}
\ifPDFTeX
  \usepackage[T1]{fontenc}
  \usepackage[utf8]{inputenc}
  \usepackage{textcomp} % provide euro and other symbols
\else % if luatex or xetex
  \usepackage{unicode-math} % this also loads fontspec
  \defaultfontfeatures{Scale=MatchLowercase}
  \defaultfontfeatures[\rmfamily]{Ligatures=TeX,Scale=1}
\fi
\usepackage{lmodern}
\ifPDFTeX\else
\setmainfont{DejaVu Serif}
\setsansfont{DejaVu Sans}
\setmonofont{DejaVu Sans Mono}
\setmathfont{STIX Two Math}
\fi
\ifPDFTeX\else
  % xetex/luatex font selection
\fi
% Use upquote if available, for straight quotes in verbatim environments
\IfFileExists{upquote.sty}{\usepackage{upquote}}{}
\IfFileExists{microtype.sty}{% use microtype if available
  \usepackage[]{microtype}
  \UseMicrotypeSet[protrusion]{basicmath} % disable protrusion for tt fonts
}{}
\makeatletter
\@ifundefined{KOMAClassName}{% if non-KOMA class
  \IfFileExists{parskip.sty}{%
    \usepackage{parskip}
  }{% else
    \setlength{\parindent}{0pt}
    \setlength{\parskip}{6pt plus 2pt minus 1pt}}
}{% if KOMA class
  \KOMAoptions{parskip=half}}
\makeatother
\usepackage{graphicx}
\graphicspath{{/Users/kalinicenkovitalijmikolajovic/.codex/worktrees/5c15/Script_for_translate_book/Output_Final/}}
\setkeys{Gin}{width=0.82\linewidth,height=0.82\textheight,keepaspectratio}
\makeatletter
\newsavebox\pandoc@box
\newcommand*\pandocbounded[1]{% scales image to fit in text height/width
  \sbox\pandoc@box{#1}%
  \Gscale@div\@tempa{\textheight}{\dimexpr\ht\pandoc@box+\dp\pandoc@box\relax}%
  \Gscale@div\@tempb{\linewidth}{\wd\pandoc@box}%
  \ifdim\@tempb\p@<\@tempa\p@\let\@tempa\@tempb\fi% select the smaller of both
  \ifdim\@tempa\p@<\p@\scalebox{\@tempa}{\usebox\pandoc@box}%
  \else\usebox{\pandoc@box}%
  \fi%
}
% Set default figure placement to htbp
\def\fps@figure{htbp}
\makeatother
\ifLuaTeX
\usepackage[bidi=basic,shorthands=off,provide=*]{babel}
\else
\usepackage[bidi=default,shorthands=off,provide=*]{babel}
\fi
\ifLuaTeX
  \usepackage{selnolig} % disable illegal ligatures
\fi
\setlength{\emergencystretch}{3em} % prevent overfull lines
\providecommand{\tightlist}{%
  \setlength{\itemsep}{0pt}\setlength{\parskip}{0pt}}
\usepackage{bookmark}
\IfFileExists{xurl.sty}{\usepackage{xurl}}{} % add URL line breaks if available
\urlstyle{same}
\hypersetup{
  pdflang={uk},
  hidelinks,
  pdfcreator={LaTeX via pandoc}}

\author{}
\date{}

\begin{document}

\pandocbounded{\includegraphics[keepaspectratio]{images/_page_315_Figure_7.jpeg}}

Рис. 42 2. Обчислення інтегралу методом Монте-Карло.

Очевидно, якщо випадковим чином кинути близько 100 дротиків у
прямокутник R, що містить графік функції, можна отримати оцінку значення
інтеграла, помноживши площу прямокутника R на відносну кількість
дротиків під графіком функції. Отже, результат нашої гри

т. е. получается УЧП с переменными коэффициентами.

\begin{verbatim}
3. Оскільки хвильове рівняння  u_{tt} = \alpha^2 u_{xx}  містить похідну другого порядку часу, щоб отримати єдине розв'язання в точці t>0, необхідно вказати  \partial sa  початкові умови:
\end{verbatim}

--- це інтегральний бал 1.

Щоб виконати ці обчислення на комп'ютері, нам потрібно якимось чином
отримати послідовність випадкових точок (подивимось, як це буде зроблено
пізніше), тобто наказати комп'ютеру «кидати дротики». Блок-схема на рис.
42.3 показує, як комп'ютер вирішить цю проблему.

\pandocbounded{\includegraphics[keepaspectratio]{images/_page_316_Figure_2.jpeg}}

\textbf{РИС.} 42.3. Блок-схема розрахунку інтегрального
\(\int_{a}^{b} f(x) dx\) методом Монте-Карло (100 кидків);
\(M = \max_{a} f(x)\) .

\section{Випадкові
числа}\label{ux432ux438ux43fux430ux434ux43aux43eux432ux456-ux447ux438ux441ux43bux430}

Перш ніж перейти до рівняння Монте-Карло з частковими похідними,
зупинимося на важливому питанні випадкових чисел. Щойно, обчислюючи
інтеграл, нам потрібно було побудувати послідовність випадкових точок
\(P_i = (x_i, y_i)\) всередині прямокутника R. Іншими словами,
x-координата точки P має бути випадковим числом з інтервалу {[}a, b{]},
а \(y_i\) від інтервалу {[}0, M{]}. Щоб знайти випадкові числа з цих
інтервалів, звертаємося до послідовності випадкових чисел, рівномірно
розподілених по сегменту {[}0, 1{]}. Тоді числа, рівномірно розподілені
на сегменті {[}a, b{]}, можна обчислити за формулою

(початкове зміщення струни), \(u_t\)(x, 0) = \(g(x)\) (початкова
швидкість струни).

Таким чином, хвильове рівняння відрізняється від рівняння
теплопровідності, для якого потрібно встановити лише одну початкову
умову.

\begin{enumerate}
\def\labelenumi{\arabic{enumi}.}
\setcounter{enumi}{3}
\tightlist
\item
  За допомогою хвильового рівняння можна описати розподіл електричного
  струму в дроті. З законів Кірхгофа отримуємо таку систему двох рівнянь
  у похідних похідних першого порядку:
\end{enumerate}

(16.3)

Тепер, звичайно, виникає питання: як отримати послідовність випадкових
чисел, \(\{r_i: i=1,2,\ldots\}\) рівномірно розподілених на {[}0,1{]}?
Найпоширенішим методом сьогодні є метод дедукцій (метод порівнянь).
Використовуючи цей метод, отримуємо послідовність випадкових цілих чисел
(наприклад, 2120, 1401, 177, 3013,\ldots), потім ліворуч від кожного
числа розміщується десяткова крапка, так що всі числа знаходяться між
нулем і одним \((0,212; 0,1401; 0,0177; 0,3013; \ldots)\) .

Отже, щоб отримати послідовність випадкових цілих чисел (між 0 і P), ми
використаємо метод віднімання.

\section{Отримання випадкових чисел за допомогою методу
відрахування.}\label{ux43eux442ux440ux438ux43cux430ux43dux43dux44f-ux432ux438ux43fux430ux434ux43aux43eux432ux438ux445-ux447ux438ux441ux435ux43b-ux437ux430-ux434ux43eux43fux43eux43cux43eux433ux43eux44e-ux43cux435ux442ux43eux434ux443-ux432ux456ux434ux440ux430ux445ux443ux432ux430ux43dux43dux44f.}

\begin{enumerate}
\def\labelenumi{\arabic{enumi}.}
\item
  Як перше випадкове ціле число, обирайте будь-яке число, що знаходиться
  між 0 і P (число P обирається наперед).
\item
  Помножте це випадкове число на певне фіксоване число.
\end{enumerate}

Попередньо обраний множник M.

\begin{enumerate}
\def\labelenumi{\arabic{enumi}.}
\setcounter{enumi}{2}
\item
  Додайте до добутку деяке фіксоване ціле число \emph{K}, обране
  заздалегідь.
\item
  Поділити отриману суму на P і залишок
\end{enumerate}

Виберіть нове випадкове число.

Метод виведення можна записати у формі формули
\(r_{i+1} \equiv (Mr_i + K) \bmod P\) , що означає, що для отримання
наступного \(r_{i+1}\) потрібно взяти попередній \(r_i\) , помножити
його на M, додати результат на K, поділити суму на P і взяти решту
поділу. Наприклад, якщо P = 100, M = 37, K = 16, \(r_0 = 15\) , то
отримується наступна послідовність випадкових чисел:

где

ж --- координата вдоль провода,

\emph{t} --- время,

\begin{verbatim}
i(x, t) — распределение тока вдоль провода,  v\left(x,\ t\right)  — распределение потенциала вдоль провода,
\end{verbatim}

С -- емкость провода на единицу длины,

G --- утечка на единицу длины,

Р --- сопротивление на единицу длины,

-- индуктивность провода на единицу длины.

Рівняння (16.3) зазвичай називають системою телеграфних рівнянь. У такій
формі вони зазвичай не використовуються. Щоб перетворити цю систему,
диференціюйте перше рівняння на x, друге на t, помножте обидві частини
на C і відніміть від першого. У результаті ми отримуємо

\[
Mr_1 + K = 2643
\]

Воспользуемся теперь вторым уравнением из (16.3)

\[
r = 71
\]

окончательно получаем

(16.4)

Y

Це рівняння для струму, відоме як телеграфне рівняння, є гіперболічним
рівнянням другого порядку (якщо, звісно, C і L не дорівнюють нулю,
інакше рівняння стає параболічним).

Напряжение описывается точно таким же уравнением:

(16.5)

\[
I = \int_{0}^{1} \int_{0}^{1} \int_{0}^{1} \int_{0}^{1} e^{-(x^{2}+y^{2}+z^{2}+\omega^{2})} dx dy dz dw,
\]

Если G = R = 0, то уравнения (16.4)---(16.5) упрощаются:

(16.6)

\[
I = \int_{0}^{1} e^{\sin x} dx.
\]

\section{\texorpdfstring{ЗАВДАННЯ - 1. Отримайте рівняння (16.5) для
\emph{and} з системи з двох рівнянь першого порядку
(16.3)}{ЗАВДАННЯ - 1. Отримайте рівняння (16.5) для and з системи з двох рівнянь першого порядку (16.3)}}\label{ux437ux430ux432ux434ux430ux43dux43dux44f---1.-ux43eux442ux440ux438ux43cux430ux439ux442ux435-ux440ux456ux432ux43dux44fux43dux43dux44f-16.5-ux434ux43bux44f-and-ux437-ux441ux438ux441ux442ux435ux43cux438-ux437-ux434ux432ux43eux445-ux440ux456ux432ux43dux44fux43dux44c-ux43fux435ux440ux448ux43eux433ux43e-ux43fux43eux440ux44fux434ux43aux443-16.3}

\begin{itemize}
\tightlist
\item
  \begin{enumerate}
  \def\labelenumi{\arabic{enumi}.}
  \setcounter{enumi}{1}
  \tightlist
  \item
    Знаючи фізичне значення кожного члена хвильового рівняння, що ви
    можете сказати про поведінку в момент розв'язання наступної задачі:
  \end{enumerate}
\end{itemize}

(YYII)

\[
r_{i+1} \equiv (3r_i + 4) \mod 7, \quad r_0 = 0.
,  0 < x < 1  ,  0 < t < \infty  ,
\]

(ГУ)

\[
I = \int_{0}^{1} \int_{0}^{1} \int_{0}^{1} e^{-(x^{2}+y^{2}+z^{2})} dx dy dz.
u(x,y) = g(x,y) =
\begin{cases} u(0, t) = 0, & 0 < t < \infty, \\ u(1, t) = 0, & 0 < t < \infty, \end{cases}
\]

Після того, як цю конкретну задачу буде розв'язано методом Монте-Карло,
ми розглянемо, як розв'язувати більш загальні задачі за допомогою цього
методу.

Щоб проілюструвати метод Монте-Карло, розглянемо гру під назвою
«Блукаючий п'яниця». Щоб грати в цю гру, потрібна таблиця з сіткою
ліній, як показано на рис. 43.1.

Перейдем к описанию правил игры.

\section{Как играть в «Блуждающего
пьяницу»?}\label{ux43aux430ux43a-ux438ux433ux440ux430ux442ux44c-ux432-ux431ux43bux443ux436ux434ux430ux44eux449ux435ux433ux43e-ux43fux44cux44fux43dux438ux446ux443}

КРОК 1. Блукання п'яниці починаються з довільної точки сітки (у нашому
випадку --- точки А).

\pandocbounded{\includegraphics[keepaspectratio]{images/_page_321_Figure_4.jpeg}}

\begin{itemize}
\item
  Початкова точка - - внутрішні точки сітки - □ - Кінцева точка p;
\item
  \(g_i\) винагороди в m. \(p_I\) Фіг. 43.1. Ігрове поле для гри
  «Блукаючий п'яниця».
\end{itemize}

КРОК 2. У кожному ході гри п'яний випадково переходить до однієї з
чотирьох суміжних точок сітки (у нашому випадку сусідні точки для A ---
це точки B, C, D і E, і ймовірність досягти кожної з цих точок становить
0,25).

Таблица 43.1

\begin{verbatim}
Вероятность окончания случайного блуждания в граничной точке p_i и величина вознаграждения g_i | Граничная точка р <sub>і</sub>                        | P_{\mathcal{A}}\left(p_{i}\right) — относительная доля блужданий, закончившихся в точке p_{i} | g_i — величина вознаграждения при достижении точки p_i |  |  |
\end{verbatim}

\textbar-------------------------------------------------------\textbar---------------------------------------------------------------------------------------------------\textbar------------------------------------------------------------\textbar--\textbar--\textbar{}

1234567891011 \textbar{}
0,040,150,030,060,170,050,060,150,030,060,160,04 \textbar{} 110000000
\textbar{} \textbar{} \textbar{}

ti

КРОК 3. Після переходу до наступної точки процес блукання відновлюється.
Таким чином, п'яниця ходить від точки до точки, поки випадково не
опиняється на межі \(p_t\) . Тут це зупиняється, і ми записуємо номер
точки \(p_t\) . Це кінець однієї невимушеної прогулянки.

КРОК 4. Повторимо 1-3 рази достатньо і знайдемо відносну кількість
прибуттів п'яниця в кожній з меж \(p_i\) . Таблиця 1. 43.1 показує
типові результати, отримані внаслідок 100 випадкових прогулянок.

\begin{verbatim}
КРОК B: Припустимо, що п'яниця отримує винагороду g_i = g\left(p_i\right) (значення граничної умови в точці p_i , якщо він досягає точки p_i ) і припустимо, що мета гри — обчислити середній R\left(A\right) винагороди за всі випадкові прогулянки, починаючи з точки A. Цей середній виграш визначається формулою
\end{verbatim}

Игра завершается, когда величина \(R\left(A\right)\) найдена. В
соответствии с данными табл. 43.1 для \(R\left(A\right)\) получаем

Теперь объясним, зачем нам понадобилось играть в такую игру.

\section{Для чего играть в «Блуждающего
пьяницу»?}\label{ux434ux43bux44f-ux447ux435ux433ux43e-ux438ux433ux440ux430ux442ux44c-ux432-ux431ux43bux443ux436ux434ux430ux44eux449ux435ux433ux43e-ux43fux44cux44fux43dux438ux446ux443}

Виявляється, що середня винагорода, про яку щойно йдеться, є приблизним
розв'язком проблеми Діріхле в точці A. Це цікаве спостереження базується
на двох фактах.

\begin{enumerate}
\def\labelenumi{\arabic{enumi}.}
\item
  Припустимо, що п'яниця почав свій шлях з точки, що лежить на Єрениці.
  Кожна така прогулянка одразу закінчується в одному й тому ж місці, і
  птах отримує винагороду \(g_i\) . Отже, середня винагорода за кожну
  граничну точку дорівнює \(g_i\) .
\item
  Тепер припустимо, що прогулянка починається з внутрішньої точки. Тоді
  стає зрозуміло, що середня винагорода за бал \(R(A)\) буде
  арифметичним середнім середніх нагород за чотири сусідні бали
\end{enumerate}

\[
\begin{cases} u(x, 0) = \sin(\pi x), & 0 \le x \le 1, \\ u_t(x, 0) = 0. & 0 \le x \le 1, \end{cases}
\]

Давайте знову запитаємо себе: чому середнє значення винагороди \(R(A)\)
наближається до розв'язання задачі Діріхле в точці A? Ми бачили, що
значення \(R(A)\) задовольняє двох

уравнениям

(во внутренних точках), \(R(A) = g_i\) (в граничных точках).

Если \(g_i\) ---это значения функции \(g(x, y)\) из граничного условия в
граничных точках \(p_i\) , то два наших уравнения точно совпадают с
двумя уравнениями, которые были получены при решении задачи Дирихле
методом конечных разностей. То есть величина \(R(A)\) соответствует
величине \(u_{i,f}\) в разностных уравнениях

(i, j)---внутренняя точка, \(u_{ij} = g_{ij}, g_{ij}\)---значение
решения в граничной точке (i, j).

Отже, значення \(R(A)\) дійсно наближає розв'язок рівняння часткових у
точці A.

Тепер можна вважати, що гра «Блукаючий п'яниця» складається з трьох
кроків.

\section{Розв'язання рівняння Лапласа за методом
Монте-Карло}\label{ux440ux43eux437ux432ux44fux437ux430ux43dux43dux44f-ux440ux456ux432ux43dux44fux43dux43dux44f-ux43bux430ux43fux43bux430ux441ux430-ux437ux430-ux43cux435ux442ux43eux434ux43eux43c-ux43cux43eux43dux442ux435-ux43aux430ux440ux43bux43e}

\begin{itemize}
\item
  Наступні правила дозволяють отримати розв'язок в одній внутрішній
  точці квадрата.
\item
  КРОК 1. Давайте зробимо кілька випадкових прогулянок, починаючи з
  точки A і закінчуючи однією з прикордонних точок. Давайте
  відстежувати, скільки разів прогулянки закінчувалися на кожній межі.
\item
  КРОК 2. Після завершення всіх прогулянок для кожної граничної точки ми
  обчислюємо відносну кількість прогулянок, які закінчилися в цій точці.
  Позначимо ці значення \(P_A\)(\(p_i\)) .
\end{itemize}

\begin{verbatim}
- КРОК 3. Обчислюємо наближений розв'язок u(A) за формулою u(A) = g_1 P_A(p_1) + g_2 P_A(p_2) + \dots + g_N P_A(p_N) ,
\end{verbatim}

де \(g_i\) --- значення функції g у точці \(p_i\) , а N --- кількість
граничних точок.

Блукаючий п'яниця можна покращити, щоб розв'язати складнішу задачу.
Нижче наведено приклад такого оновлення.

Розв'язок задачі Діріхле з змінними коефіцієнтами.

Рассмотрим эллиптическую краевую вадачу в квадрате

Дайте физическую интерпретацию этой задачи.

\begin{enumerate}
\def\labelenumi{\arabic{enumi}.}
\setcounter{enumi}{3}
\tightlist
\item
  Для уравнения \(u_{tt} = u_{xx}\) найдите все решения вида
\end{enumerate}

\begin{verbatim}
u_{xx} + (\sin x) u_{yy} = 0 , 0 < x < \pi , 0 < y < \pi , (\Gamma Y) u(x, y) = g(x, y) на границе квадрата.
\end{verbatim}

\begin{verbatim}
Щоб розв'язати цю задачу, замінимо u_{xx} , u_{yy} та \sin x наступним чином:
\end{verbatim}

Будет ли решением сумма двух решений такого вида?

\section{ФОРМУЛА ДАЛАМБЕРА МЕТА ЛЕКЦІЇ: Построить решение задачи Коши
для волнового
уравнения}\label{ux444ux43eux440ux43cux443ux43bux430-ux434ux430ux43bux430ux43cux431ux435ux440ux430-ux43cux435ux442ux430-ux43bux435ux43aux446ux456ux457-ux43fux43eux441ux442ux440ux43eux438ux442ux44c-ux440ux435ux448ux435ux43dux438ux435-ux437ux430ux434ux430ux447ux438-ux43aux43eux448ux438-ux434ux43bux44f-ux432ux43eux43bux43dux43eux432ux43eux433ux43e-ux443ux440ux430ux432ux43dux435ux43dux438ux44f}

и \emph{\{yy\}=\(u_{l+1,\ j}\)-2u}\{l,~j\}+\(u_{l-1,\ j}\), (центральные
разностные производные),

і підставляють їх у рівняння похідних похідних. Відповідна сітка
показана на рис. 43.2. Розв'язавши отримане рівняння за \(u_{i,l}\) ,
отримуємо

Посмотрим внимательно на это уравнение. Коэффициенты при
\(u_{i+1,\ j},\ u_i,\ j+i,\ u_l,\ j-1\) н \(u_{i-1},\ j\) положительны,
и их сумма равна единице. Другими словами, решение \(u_i,\ j\) является
взвешенным средним решений в четырех соседних точках. Следовательно,
можно Р И С

\pandocbounded{\includegraphics[keepaspectratio]{images/_page_324_Picture_9.jpeg}}

. 43.2. Узлы сетки для случайного маршруга.

модифицировать «Блуждающего пьяницу» так, чтобы вероятности перехода в
соседние точки были не 0,25, а равнялись бы коэффициенту в
соответствующем члене. Другими словами, если пьяница находится в точке
(i, j), то он переходит в точку

Ця задача описує рух нескінченної струни за заданих початкових умов.
Його розв'язав у 1750 році французький математик д'Аламбер. Формула
Д'Аламбера

с вероятностью \(\frac{1}{2(1+\sin x_j)}\) , \((i, j-1)\) с вероятностью
\(\frac{1}{2(1+\sin x_j)}\) , \((i+1, j)\) с вероятностью
\(\frac{\sin x_j}{2(1+\sin x_j)}\) , \((i-1, j)\) с вероятностью
\(\frac{\sin x_j}{2(1+\sin x_j)}\) .

В остальном игра не меняется. Модификация для других задач может
оказаться более хитроумной, но основная идея остается той же. Читатель
может заняться разработкой своих игр для решения других задач.
Параболический случай рассмотрен в разделе «Задачи». \# ЗАУВАЖЕННЯ 1.
Отметим, что если относительные частоты попадания \(P_A(p_l)\) известны,
то можно легко найти величину \(u\left(A\right)\) при других граничных
условиях \(g_l\) . Для этого достаточно подставить \(g_l\) в формулу

полегшує пошук розв'язку, якщо початкові умови відомі. Крім того, це
дозволяє дати цікаву фізичну інтерпретацію розв'язку мовою хвиль, що
рухаються.

\begin{verbatim}
Як читач, ймовірно, пам'ятає, у параболічному випадку ми спочатку розв'язали задачі дифузії на скінченному відрізку (методом розділення змінних), а лише потім переходили до необмежених інтервалів (-  \infty < x < \infty  ), де застосовували перетворення Фур'є. У гіперболічному випадку ми зробимо навпаки. Почнемо з розв'язання одномірного хвильового рівняння на всій дійсній осі, тобто з задачі Коші
\end{verbatim}

(17.1)

а не разыгрывать вновь случайные прогулки.

\begin{enumerate}
\def\labelenumi{\arabic{enumi}.}
\setcounter{enumi}{1}
\tightlist
\item
  Во многих случаях исследователя интересует решение уравнения с
  частными производными в одной точке. Если граница области сложна и
  уравнение зависит от трех или четырех независимых переменных, метод
  Монте-Карло может оказаться удобным. Методы Монте-Карло возникли при
  решении очень трудных задач диффузии нейтронов. Эти задачи невозможно
  решать аналитически. \# ЗАДАЧІ
\item
  Построить блок-схему игры «Блуждающий пьяница» для решения задачи
\end{enumerate}

(УЧП)

, 0 \textless{} x \textless{} 1 , 0 \textless{} y \textless{} 1 , (ГУ)
\(u(x, y)\) = \(g(x, y)\) на границе,

\begin{itemize}
\item
  у вузлах внутрішньої сітки, якщо кількість вузлів сітки довільна. 2.
  Напишіть програму для розв'язання задачі відповідно до блок-схеми з
  завдання 1.
\item
  \begin{enumerate}
  \def\labelenumi{\arabic{enumi}.}
  \setcounter{enumi}{2}
  \tightlist
  \item
    Як модифікувати гру для вирішення проблеми
  \end{enumerate}
\end{itemize}

(УЧП)

содержащей только начальные условия. Эту задачу можно было бы решить
либо с помощью преобразования Фурье (по переменной t), либо с помощью
преобразования Лапласа (по переменной t). Однако мы воспользуемся другим
методом (методом канонических координат), который, как мы надеемся,
заинтересует читателя. Этот метод базируется на тех же идеях, что и
метод перехода к движущейся системе отсчета, с которым мы познакомились
в лекции 15. Итак, приступим к решению задачи (17.1).

\section{Решение одномерного волнового уравнения. Формула
Даламбера}\label{ux440ux435ux448ux435ux43dux438ux435-ux43eux434ux43dux43eux43cux435ux440ux43dux43eux433ux43e-ux432ux43eux43bux43dux43eux432ux43eux433ux43e-ux443ux440ux430ux432ux43dux435ux43dux438ux44f.-ux444ux43eux440ux43cux443ux43bux430-ux434ux430ux43bux430ux43cux431ux435ux440ux430}

Решение задачи (17.1) разобьем на несколько шагов:

ШАГ 1 (Замена координат (x, t) новыми каноническими координатами
\((\xi, \eta)\) ).

Для решения задачи (17.1) воспользуемся тем, что если заменить две
независимые переменные x и t новыми пространственновременными
координатами

, 0 \textless{} x \textless{} 1 , 0 \textless{} y \textless{} 1 , (ГУ)
\(u(x, y)\) = \(g(x, y)\) на границе?

Как выглядит в этой игре случайная прогулка?

\begin{enumerate}
\def\labelenumi{\arabic{enumi}.}
\setcounter{enumi}{3}
\tightlist
\item
  Можете ли вы так модифицировать игру, чтобы решить задачу
\end{enumerate}

\[
\begin{array}{l} u_{xx} + u_{yy} = 0, \quad 0 < x < 1, \quad 0 < y < 1, \\ u(x, 1) = 0, \quad 0 < x < 1, \quad 0 < y < 1, \\ u(x, 0) = 0, \quad 0 < x < 1, \quad 0 < y < 1, \\ u(0, y) = 1, \quad \frac{\partial u}{\partial x}(1, y) = 0 \end{array}
\]

то уравнение

\begin{verbatim}
\begin{aligned}
(\text{УЧП})\quad & u_t = \alpha^2 u_{xx} - \beta u_x - \gamma u, && 0 < x < 1,\ 0 < t < \infty, \\
(\text{ГУ})\quad & u(0,t) = f(t),\quad u(1,t) = g(t), && 0 < t < \infty, \\
(\text{НУ})\quad & u(x,0) = \varphi(x), && 0 \leqslant x \leqslant 1.
\end{aligned}
\end{verbatim}

преобразуется в уравнение

ПРИМІТКА. Замінимо це рівняння на скінченну різницю Кренка--Ніколсона,
виразимо \(u_{i+1,j}\) через п'ять значень \(u_{i+1,j-1}\) ,
\(u_{i+1,j+1}\) , \(u_{i,j+1}\) , \(u_{i,j+1}\) , \(u_{i,j+1}\) у
суміжних точках.

\section{КАЛЬКУЛЮС ВАРІАЦІЙ (РІВНЯННЯ
ЕЙЛЕРА-ЛАГРАНЖА)}\label{ux43aux430ux43bux44cux43aux443ux43bux44eux441-ux432ux430ux440ux456ux430ux446ux456ux439-ux440ux456ux432ux43dux44fux43dux43dux44f-ux435ux439ux43bux435ux440ux430-ux43bux430ux433ux440ux430ux43dux436ux430}

МЕТА ЛЕКЦІЇ: Ввести поняття функціоналу, тобто функції від функції, і
пояснити, як функціонали природно виникають у фізиці. Типовим
функціоналом є інтеграл

\[
J[y] = \int_a^b F(x, y, y')\,dx,
\]

який є функцією від функції \(y\) (при цьому вважається, що
підінтегральна функція \(F(x, y, y')\) задана). Приклад функціоналу:

\[
J[y] = \int_0^1 \left[y^2(x) + y'^2(x)\right] dx.
\]

Ми покажемо, як знайти функцію \(\overline{y}(x)\), що мінімізує
функціонал \(J[y]\). Виявляється, що мінімізуюча функція
\(\overline{y}\) повинна задовольняти так зване рівняння
Ейлера-Лагранжа. Воно відіграє ту саму роль, що й необхідна умова
мінімуму

\[
\frac{df(x)}{dx} = 0
\]

для функції \(f(x)\) у точці \(x\) в диференціальному численні.

Варіаційне числення тісно пов'язане з диференціальними рівняннями, але,
на жаль, більшість студентів його не вивчає. У цій лекції (і наступній)
подається вступ до варіаційного числення та показано, як можна
розв'язувати диференціальні рівняння в частинних похідних на основі
варіаційних принципів.

Варіаційне числення виникло одночасно з математичним аналізом у зв'язку
з розв'язанням задач максимізації та мінімізації функцій від функцій,
тобто функціоналів. Першою задачею варіаційного числення була задача про
брахістохрону, сформульована Йоганном Бернуллі у 1696 році. У цій задачі
потрібно знайти криву \(y(x)\) так, щоб мінімізувати час спуску без
тертя з однієї точки в іншу (див. рис. 44.1).

Бернуллі показав, що час спуску записується у вигляді

\[
T = \int_0^T dt = \int_a^L \frac{dt}{ds}\,ds = \int_0^L \frac{ds}{v}
= \frac{1}{\sqrt{2mg}} \int_0^L \frac{ds}{\sqrt{y}}
= \frac{1}{\sqrt{2mg}} \int_a^b \sqrt{\frac{1+y'^2}{y}}\,dy
\]

і, отже, є функціоналом \(T\{y\}\) від функції. Оскільки багато
функціоналів мають подібну будову, ми зосередимося на вивченні
функціоналу загального вигляду

(44.1) \[
J[y] = \int_a^b F(x, y, y')\,dx.
\]

Тепер сформулюємо головну мету лекції: знайти функцію, яка доставляє
мінімум (або максимум) функціоналу (44.1). Стратегія пошуку буде такою
ж, як і при пошуку мінімуму функції у звичайному диференціальному
численні. Там ми шукали критичні точки з умови \(f'(x)=0\). У
варіаційному численні все дещо складніше, бо аргументом є не число, а
функція. Проте загальний підхід той самий: ми обчислимо функціональну
похідну за функцією \(y(x)\) і прирівняємо її до нуля. Нове рівняння
буде аналогом умови \(df(x)/dx = 0\), але тепер це буде звичайне
диференціальне рівняння, відоме як рівняння Ейлера-Лагранжа.

Мінімізація функціоналу

\[
J[y] = \int_a^b F(x, y, y')\,dx
\]

Розглянемо задачу знаходження функції \(y(x)\), яка мінімізує цей
функціонал у класі гладких функцій, що задовольняють граничним умовам

\[
y(a) = A, \qquad y(b) = B.
\]

(див. рис. 44.2).

Нехай шукана функція \(y = \overline{y}(x)\) існує, і розглянемо малу
варіацію цієї функції, тобто функцію
\(\overline{y} + \varepsilon\eta(x)\), де \(\varepsilon\) --- мале
число, а \(\eta(x)\) --- гладка функція, яка задовольняє граничним
умовам \(\eta(a)=\eta(b)=0\).

Тоді, якщо обчислити інтеграл \(J\) для близької функції
\(\overline{y} + \varepsilon\eta\), функціонал зросте, тобто

\[
J[\overline{y}] \leqslant J[\overline{y} + \varepsilon\eta]
\]

для всіх \(\varepsilon\). Іншими словами, графік
\(\varphi(\varepsilon)=J[\overline{y}+\varepsilon\eta]\) як функції
\(\varepsilon\) має вигляд, показаний на рис. 44.3.

\pandocbounded{\includegraphics[keepaspectratio]{images/_page_329_Figure_5.jpeg}}

Рис. 44.1. Задача про брахістохрону (з цієї задачі почалося варіаційне
числення).

\pandocbounded{\includegraphics[keepaspectratio]{images/_page_329_Figure_7.jpeg}}

Рис. 44.2. Варіація функції: a --- варіація функції
\(\overline{y}(x)+\varepsilon\eta(x)\); б --- мінімізуюча функція
\(\overline{y}(x)\).

З рис. 44.3 видно, що тут потрібно обчислити похідну від

\[
\varphi(\varepsilon) = J[\overline{y} + \varepsilon\eta]
\]

за \(\varepsilon\), покласти \(\varepsilon = 0\) і прирівняти отриманий
вираз до нуля, тобто

\[
\frac{d\varphi(\varepsilon)}{d\varepsilon}
= \frac{d}{d\varepsilon} J[\overline{y} + \varepsilon\eta]\bigg|_{\varepsilon=0}
= \int_a^b \left[
\frac{\partial F}{\partial \overline{y}}\eta(x)
+ \frac{\partial F}{\partial \overline{y'}}\eta'(x)
\right] dx = 0.
\]

(Читач має виконати це обчислення самостійно.) Інтегруючи частинами,
дістаємо

\[
\frac{d\varphi(\varepsilon)}{d\varepsilon}
\equiv \int_a^b \left\{
\frac{\partial F}{\partial \overline{y}} - \frac{d}{dx}\left[\frac{\partial F}{\partial \overline{y'}}\right]
\right\} \eta(x)\,dx = 0.
\]

Оскільки цей інтеграл дорівнює нулю для будь-якої функції \(\eta(x)\),
яка задовольняє граничним умовам \(\eta(a)=\eta(b)=0\), то
підінтегральний вираз має дорівнювати нулю, тобто

\pandocbounded{\includegraphics[keepaspectratio]{images/_page_330_Figure_6.jpeg}}

Рис. 44.3. Графік функції \(J[y+\varepsilon\eta]\) в околі
\(\varepsilon=0\).

(44.2) \[
\frac{\partial F}{\partial \overline{y}} - \frac{d}{dx}\left[\frac{\partial F}{\partial \overline{y'}}\right] = 0
\]

(рівняння Ейлера-Лагранжа).

Рівняння (44.2) називається рівнянням Ейлера-Лагранжа. Хоча в загальному
вигляді воно здається складним, після підстановки конкретної функції
\(F(x,y,y')\) воно перетворюється на звичайне диференціальне рівняння
другого порядку відносно невідомої функції \(\overline{y}(x)\). Отже,
щоб визначити мінімізуючу функцію \(\overline{y}\), треба розв'язати
рівняння Ейлера-Лагранжа.

Отже, ми показали, що якщо функція \(y(x)\) мінімізує функціонал
\(J[y] = \int_a^b F(x,y,y')\,dx\) у класі гладких функцій із граничними
умовами \(y(a)=A\) і \(y(b)=B\), то вона повинна задовольняти рівнянню

\[
\frac{\partial F}{\partial y} - \frac{d}{dx}\left[\frac{\partial F}{\partial y'}\right] = 0.
\]

(Для спрощення позначень ми прибрали риску над \(\overline{y}\).) Для
ілюстрації теорії розглянемо приклад.

Пошук мінімальної функціональності
\(J[y] = \int_{0}^{1} [y^2 + y'^2] dx\) .

Спробуємо знайти функцію \(y(x)\), яка проходить через точки (0,0) і
(1,1) і мінімізує функціональність

\[
P_{0}
\]

З контексту зрозуміло, що шукана функція має бути диференційовною,
оскільки інтегративний вираз залежить від y'. Щоб визначити
\(\bar{y}(x)\) , запишемо рівняння Ейлера--Лагранжа (разом із граничною
умовою \(y(0)\)=0 та \(y(1)\)=1).

\begin{verbatim}
,  (x,  y, z) \in R^3  ,  u = 0   u_t =
\end{verbatim}

\(y(0) = 0,\) \(y(1) = 1.\) Але \(F(x, y, y') = y^2 + y'^2\) і, отже,

\begin{enumerate}
\def\labelenumi{\arabic{enumi}.}
\setcounter{enumi}{3}
\tightlist
\item
  Решите аналогичную (см. задачу 3) задачу для двумерного уравнения
\end{enumerate}

(УЧП)

(Продифференцировали \(F\) по \(y\) и \(y'.\) )

Уравнение Эйлера --- Лагранжа принимает вид

Зверніть увагу, що в перших двох задачах коефіцієнти є сталими. Читач
повинен розуміти, що параметр b має бути достатньо малим, а необережний
ряд може відрізнятися

\section{ЗАДАЧИ}\label{ux437ux430ux434ux430ux447ux438}

\begin{verbatim}
- 1. Подставьте разложение (46.8) в задачу (46.7) и получите последовательность задач P_0,\ P_1,\ P_2,\ \dots . 2. Покажите, что нелинейную задачу
\end{verbatim}

откуда

(НУ)

Решая это простое дифференциальное уравнение с граничными условиями
\(y(0)\)=0 и \(y(1)\)=1, получаем

\[
u(r, \theta) = u_0(r, \theta) + \frac{1}{4}u_1(r, \theta)
\]

.

График этой функции изображен на рис. 44.4. Ясно, что любая другая
гладкая функция, удовлетворяющая тем же граничным условиям, даст большее
значение функционала \(J(y)\).

\pandocbounded{\includegraphics[keepaspectratio]{images/_page_332_Figure_5.jpeg}}

Рис. 44 4. Все допустимые гладкие кривые удовлетворяют граничным
условиям \(y(0)\) = 0 и \(y(1)\) = 1: a---допустимая кривая; \(\delta\)
---минимизирующая функция \(y(x) = 0.42e^x - 0.42e^{-x}\) . \#
ЗАУВАЖЕННЯ 1. Уравнение Эйлера---Лагранжа аналогично равенству нулю
производной в дифференциальном исчислении. Если читатель помнит, то на
самом деле равенство нулю производной не является достаточным условием
экстремума. Например, производная функции \(f(x) = x^3\) при x = 0 равна
нулю, но эта точка не является нп точкой минимума, ни точкой максимума.
Так же и с уравнением Эйлера---Лагранжа: оно является только необходимым
условием экстремума, но не достаточным. Значит, экстремум может
достигаться на других функциях. Очень часто, однако, решение уравнения
Эйлера---Лагранжа по существу задачи дает локальный или даже глобальный
минимум.

\begin{enumerate}
\def\labelenumi{\arabic{enumi}.}
\setcounter{enumi}{1}
\tightlist
\item
  Если у (х) минимизирует функционал
\end{enumerate}

\[
\Delta u = 0,
\]

то площадь поверхности, образованной вращением кривой \(y(x)\) вокруг
оси x, будет минимальной (рис. 44.5). Решением уравнения Эйлера ---
Лагранжа в этом случае будет

цепная линия

\[
q_{xx} + \varphi_{uu} = 0
\]

где константы \(\alpha\) и \(\beta\) определяются из условий \(y(a)\) =
A и \(y(b)\) = B (т. е. цепная линия должна проходить через заданные
точки). Решить уравнение Эйлера---Лагранжа в данном случае достаточно
трудно, поэтому мы рекомендуем читателю только проверить, что цепная
линия ему удовлетворяет.

\begin{enumerate}
\def\labelenumi{\arabic{enumi}.}
\setcounter{enumi}{2}
\tightlist
\item
  В этой лекции мы нашли, что функция
\end{enumerate}

\[
q_{uu} + q_{vv} = 0
\]

доставляет минимум функционалу Р

\pandocbounded{\includegraphics[keepaspectratio]{images/_page_333_Picture_8.jpeg}}

ис. 44 5. Минимальная поверхность вращения.

\begin{enumerate}
\def\labelenumi{\arabic{enumi}.}
\setcounter{enumi}{3}
\tightlist
\item
  Основные законы физики чаще всего формулируются на языке вариационных
  принципов, а не дифференциальных уравнений. В качестве примеров можно
  привести принцип Ферма (свет при распространении из одной точки в
  другую выбирает путь, которому соответствует наименьшее время
  распространения) или принцип Гамильтона (в консервативном поле частица
  движется так, что интеграл действия
\end{enumerate}

\[
w = f(z)
\]

(кинетическая энергия--- потенциальная энергия) \(dt\)

будет минимальным). Таким образом, физические явления развиваются только
так, что эти функционалы принимают минимальные значения.

Основные идеи вариационного исчисления можно распространить на
функционалы, зависящие от функции нескольких независимых переменных,
вида П

\pandocbounded{\includegraphics[keepaspectratio]{images/_page_335_Picture_2.jpeg}}

ринцип Гамильтона утверждает, что материальная точка будет двигаться
так, что интеграл

\begin{verbatim}
\left\{
\end{verbatim}

будет минимален. Если этот закоп справедлив, то получите
дифференциальное уравнение движения материальной точки.

\begin{enumerate}
\def\labelenumi{\arabic{enumi}.}
\setcounter{enumi}{2}
\tightlist
\item
  Покажите, что в классе гладких функций, удовлетворяющих граничным
  условиям \(y(0)\)=0 и \(y(\pi/2)=1\) , минимум функционала
\end{enumerate}

Чи знаєте ви, як виглядає графік цієї функції з сегментом {[}0, 1{]}?
Подумай про це. Зверніть увагу, що синусоїдальне перетворення функції
\(f(x)\) визначається лише для додатних цілих чисел n.~Іншими словами,
скінченний синус і косинус перетворюють функції у чисельні
послідовності.

\pandocbounded{\includegraphics[keepaspectratio]{images/_page_343_Figure_2.jpeg}}

\begin{verbatim}
Далсе. Приближенное решение   u_n(x,y) = a_1 \varphi_1 + a_2 \varphi_2 + \dots + a_n \varphi_n  .
\end{verbatim}

В это время можно численно найти эначение функционала J{[}\(u_n\){]} .
Если необходимо найти значение решения \(u_n\)(x,y) в некоторой точке,
можно написать небольшую программу эля быполнения отой операции. При
необходимости можно повторить вычисления по этой программе с другими
значениями п Рис. 452. Блок-схема метода Ритца.

\section{РЕШЕНИЕ УРАВНЕНИЙ С ЧАСТНЫМИ ПРОИЗВОДНЫМИ МЕТОДАМИ ТЕОРИИ
ВОЗМУЩЕНИЙ МЕТА ЛЕКЦІЇ: Показать, как различные сложные задачи
(нелинейные уравнения с переменными коэффициентами, области неправильной
формы и т. д.) можно решить методом возмущения более простых задач.
Иными словами, мы хотим показать, как следует изменить решение простой
задачи, чтобы оно давало приближенное решение трудной
задачи.}\label{ux440ux435ux448ux435ux43dux438ux435-ux443ux440ux430ux432ux43dux435ux43dux438ux439-ux441-ux447ux430ux441ux442ux43dux44bux43cux438-ux43fux440ux43eux438ux437ux432ux43eux434ux43dux44bux43cux438-ux43cux435ux442ux43eux434ux430ux43cux438-ux442ux435ux43eux440ux438ux438-ux432ux43eux437ux43cux443ux449ux435ux43dux438ux439-ux43cux435ux442ux430-ux43bux435ux43aux446ux456ux457-ux43fux43eux43aux430ux437ux430ux442ux44c-ux43aux430ux43a-ux440ux430ux437ux43bux438ux447ux43dux44bux435-ux441ux43bux43eux436ux43dux44bux435-ux437ux430ux434ux430ux447ux438-ux43dux435ux43bux438ux43dux435ux439ux43dux44bux435-ux443ux440ux430ux432ux43dux435ux43dux438ux44f-ux441-ux43fux435ux440ux435ux43cux435ux43dux43dux44bux43cux438-ux43aux43eux44dux444ux444ux438ux446ux438ux435ux43dux442ux430ux43cux438-ux43eux431ux43bux430ux441ux442ux438-ux43dux435ux43fux440ux430ux432ux438ux43bux44cux43dux43eux439-ux444ux43eux440ux43cux44b-ux438-ux442.-ux434.-ux43cux43eux436ux43dux43e-ux440ux435ux448ux438ux442ux44c-ux43cux435ux442ux43eux434ux43eux43c-ux432ux43eux437ux43cux443ux449ux435ux43dux438ux44f-ux431ux43eux43bux435ux435-ux43fux440ux43eux441ux442ux44bux445-ux437ux430ux434ux430ux447.-ux438ux43dux44bux43cux438-ux441ux43bux43eux432ux430ux43cux438-ux43cux44b-ux445ux43eux442ux438ux43c-ux43fux43eux43aux430ux437ux430ux442ux44c-ux43aux430ux43a-ux441ux43bux435ux434ux443ux435ux442-ux438ux437ux43cux435ux43dux438ux442ux44c-ux440ux435ux448ux435ux43dux438ux435-ux43fux440ux43eux441ux442ux43eux439-ux437ux430ux434ux430ux447ux438-ux447ux442ux43eux431ux44b-ux43eux43dux43e-ux434ux430ux432ux430ux43bux43e-ux43fux440ux438ux431ux43bux438ux436ux435ux43dux43dux43eux435-ux440ux435ux448ux435ux43dux438ux435-ux442ux440ux443ux434ux43dux43eux439-ux437ux430ux434ux430ux447ux438.}

Например, мы покажем, как нелинейную задачу

(УЧП)

\[
w = \ln \left[ \frac{z-1}{z+1} \right].
,   0 < r < 1  ,
\]

\begin{verbatim}
(ГУ)   u(1, \theta) = \cos \theta  ,   0 \le \theta \le 2\pi  .
\end{verbatim}

Очень часто задачу, решение которой не известно, можно путем
непрерывного изменения каких-то параметров свести к близкой, но легко
решаемой задаче. Тогда и решение исходной задачи будет близко решению
модифицирующей задачи.

Например, уравнение Лапласа

\[
\begin{cases} 1, & n = 1, \\ 0, & n = 2, 3, \dots, \end{cases}
\]

в помощью семейства уравнений

, \(0 \le \varepsilon \le 1\) ,

\[
\right.
\]

, \(0 < r < 1\) , \(0 \le \theta < 2\pi\) , \(u(1, \theta) = g(\theta)\)
, \(0 \le \theta < 2\pi\) ,

при одних граничных условиях (например, \(g(\theta) = \cos \theta\) )
имеет два различных вещественных решения, а при других (например,
\(g(\theta) = A \cos \theta\) с A \textgreater{} 20,65) эта задача не
имеет ни одного вещественного решения. --- Прим. ред.

4) Краевая задача

можно непрерывно модифицировать в нелинейное уравнение
\(\Delta u = u^2 = 0\) 1).

На рис. 46.1 изображена схема решения нелинейной задачи путем сведения
ее к решению уравнения Лапласа2)

\end{document}
